\chapter{因子预处理：为什么你的"新因子"只是小盘股的换皮}\label{ch:prep}

\begin{keybox}[本章要回答的问题]
去极值该用 3$\sigma$ 还是 MAD？标准化用 z-score 还是排名？中性化到底在剥离什么？
为什么"中性化后 IC 一定更高"是错的？
\end{keybox}

\section{为什么必须预处理}

原始因子值直接使用会遇到三个问题：

\begin{enumerate}
  \item \textbf{极端值主导}。A 股有大量"净资产为负 → BP 爆炸""净利润接近零 → EP 爆炸"
        的样本。一个 BP 等于 500 的股票会把线性回归的系数完全带偏。
  \item \textbf{量纲不可比}。换手率是百分数（0--30），ROE 是百分数（$-$50--50），
        市值是元（$10^8$--$10^{12}$）。不加处理无法相加。
  \item \textbf{风格暴露}。你的"新因子"可能只是"小盘股"或"医药行业"的换皮。
\end{enumerate}

\section{去极值}

\subsection{三种方法}

\begin{description}
  \item[3$\sigma$ 法] 把超过 $\mu \pm 3\sigma$ 的值截断。问题是 $\mu,\sigma$ 本身
        被极端值污染——厚尾分布下 $\sigma$ 会被拉大，导致"该截的没截掉"。
  \item[百分位法] 把超过 $1\%$ 与 $99\%$ 分位的值截断。简单稳健，但在样本很少时
        分位数估计不稳。
  \item[MAD 法] 用中位数绝对偏差：
        \begin{equation}
          \text{MAD} = \mathrm{median}\big(|x_i - \mathrm{median}(x)|\big), \qquad
          \hat\sigma = 1.4826\,\text{MAD}.
        \end{equation}
        常数 $1.4826$ 使 $\hat\sigma$ 在正态分布下与标准差一致。
        截断区间取 $[\mathrm{median} - n\hat\sigma,\ \mathrm{median} + n\hat\sigma]$，
        $n$ 常取 5。
\end{description}

\begin{keybox}[本书的选择]
截面 MAD 法，$n=5$。理由：中位数与 MAD 的崩溃点是 50\%，
即使有 20\% 的极端样本也不影响估计；而 $3\sigma$ 法的崩溃点是 0\%。
\end{keybox}

\section{标准化}

\subsection{z-score 与排名}

\begin{align}
  z_{i} &= \frac{x_i - \mathrm{mean}(x)}{\mathrm{std}(x)}, \\
  r_{i} &= \mathrm{rank}(x_i)/N - \tfrac12 \;\in\; \left(-\tfrac12, \tfrac12\right].
\end{align}

\begin{center}
\begin{tabular}{@{}lll@{}}
\toprule
 & \textbf{z-score} & \textbf{排名} \\
\midrule
抗极值 & 弱（受尾部影响） & 强 \\
保留幅度 & 是 & 否 \\
分布形状 & 近似正态 & 均匀 \\
适合 & 线性模型、组合权重 & 稳健评价、IC 计算 \\
\bottomrule
\end{tabular}
\end{center}

实践中常见的折中：\textbf{先去极值，再 z-score}。这样既保留了幅度，又不被极端值主导。

\section{中性化：剥离你不想要的风险}

\subsection{在剥离什么}

假设你构造了一个"研发投入强度"因子，回测发现 IC 很高。但仔细一看：
高研发投入的公司集中在创业板，而创业板在样本期内整体跑赢。
那么你的因子赚的其实是"创业板 beta"，而不是"研发"本身。

中性化就是把这个"其实不是你的"部分剥离掉。

\subsection{行业中性化}

最简单的做法是行业内 demean：

\begin{equation}
  \tilde x_{i} = x_i - \mathrm{mean}_{j \in \text{行业}(i)}(x_j).
\end{equation}

等价于对行业哑变量做无截距回归取残差。

\subsection{行业 + 市值中性化}

更完整的做法是截面回归：

\begin{equation}
  x_{i,t} = \alpha_t + \sum_{k=1}^{K-1}\beta_{k,t}\,\mathbb{1}[i\in \text{行业}_k]
  + \gamma_t \ln(\text{MktCap}_{i,t}) + \varepsilon_{i,t},
  \label{eq:neutral}
\end{equation}

取残差 $\varepsilon_{i,t}$ 作为中性化后的因子值。

\begin{proposition}[残差的正交性]
由最小二乘的正规方程，残差 $\hat\varepsilon$ 与设计矩阵的每一列都正交：
\begin{equation}
  \sum_i \hat\varepsilon_i = 0, \qquad
  \sum_i \hat\varepsilon_i \mathbb{1}[i\in k] = 0\ \forall k, \qquad
  \sum_i \hat\varepsilon_i \ln(\text{MktCap}_i) = 0.
\end{equation}
因此中性化后的因子在截面上\textbf{行业均值为零、与对数市值零相关}。
这正是"剥离"的数学含义。
\end{proposition}

\subsection{为什么不用分组 demean}

分组 demean 只能处理离散分组（行业），无法同时控制连续的市值。
而且它无法处理"行业内市值差异"带来的残余暴露。
回归可以一次性处理两者，且残差严格正交，是更干净的做法。

\subsection{几个工程细节}

\begin{enumerate}
  \item \textbf{市值用对数}。市值分布严重右偏，取对数后接近正态；
        直接用原始市值会让极少数巨型股主导回归。
  \item \textbf{哑变量丢一列}。$K$ 个行业要用 $K-1$ 个哑变量加截距，
        否则完全共线。
  \item \textbf{截面做，不时序做}。因子有效性是横截面问题，
        每天独立回归一次。
  \item \textbf{样本太少时放弃}。当天有效样本 $<30$ 时，回归系数不稳，宁可不中性化。
  \item \textbf{二次项通常不需要}。线性项已经足够，加二次项会增加自由度与过拟合风险；
        除非你明确怀疑存在 U 形关系且样本量足够。
\end{enumerate}

\section{一个反直觉的结论：中性化不总是提升 IC}

很多教程写"中性化后 IC 会提高"。这是错的。

\begin{itemize}
  \item 如果因子的信息部分来自行业轮动（例如"景气度向上行业的股票"），
        行业中性化会把真实 alpha 一起剥掉，IC 反而下降。
  \item 如果因子本身就是一个风格因子（例如规模因子），
        对市值中性化等于把它归零。
\end{itemize}

\begin{warnbox}[正确的判断方式]
中性化前后各算一遍 IC，看\textbf{降低幅度}。
\begin{itemize}
  \item IC 几乎不变：因子本身与行业/市值无关，中性化是"买个保险"。
  \item IC 明显下降：因子有相当一部分收益来自风格暴露，中性化后剩下的才是真 alpha。
  \item IC 反而上升：说明原始因子被风格噪声污染，中性化起了降噪作用。
\end{itemize}
三种情况都是有用信息，但只有第三种能说"中性化提升了因子"。
\end{warnbox}

\section{实现：向量化的截面回归}

逐日循环回归在小样本上够用，但在 $250$ 天 $\times$ 40 个因子上会明显变慢。
一个有效的优化是：\textbf{同一天的所有因子共用同一个设计矩阵}，
因此对 $X$ 做一次 QR 分解，就能把整块因子矩阵投影掉。

\begin{pycode}[批量残差化]
def neutralize_block(Y, X):
    """Y: (n, m) 因子矩阵；X: (n, k) 设计矩阵。返回残差 (n, m)。"""
    Q, _ = np.linalg.qr(X)        # 只分解一次
    return Y - Q @ (Q.T @ Y)      # 一次性投影掉所有因子
\end{pycode}

复杂度从 $O(m\,n k^2)$ 降到 $O(nk^2 + mnk)$，在 $m=40$、$n=4000$、$k=32$ 时
大约快 10 倍。

\section{预处理流程的完整顺序}

顺序不能乱，理由如下：

\begin{enumerate}
  \item \textbf{剔除不可交易样本}：ST、停牌、新股、流动性不足。
        先剔除，否则它们的极端值会污染去极值的统计量。
  \item \textbf{去极值}：MAD，$n=5$。
  \item \textbf{标准化}：截面 z-score。
  \item \textbf{中性化}：行业 + 对数市值回归取残差。
  \item \textbf{再次标准化}：中性化会改变量纲，需重新标准化以便加权。
\end{enumerate}

\begin{warnbox}[一个容易被忽略的顺序问题]
如果先中性化再去极值，去极值会破坏残差的正交性——
你截断掉的极端值，恰好可能是与市值高度相关的那些点。
正确顺序是\textbf{先去极值、再中性化}。
\end{warnbox}

\section{本章小结}

\begin{enumerate}
  \item 去极值用 MAD 而不是 $3\sigma$，因为后者自身被极端值污染。
  \item 标准化保留幅度去 z-score，追求稳健去排名；实践中"先去极值再 z-score"最常用。
  \item 中性化 = 截面回归取残差，残差与行业、市值严格正交。
  \item 中性化不总提升 IC；IC 下降说明你的因子有相当部分是风格暴露。
  \item 处理顺序：剔除不可交易 → 去极值 → 标准化 → 中性化 → 再标准化。
\end{enumerate}
